\section{Principle Ideal Domains}

\begin{definition}[Principle Ideal]
  Given a ring R and an element \(\mathrm{a} \in \mathrm{R}\), the principal ideal generated by a (denoted as \(\langle a\rangle\) or (a)) is defined as: 
\begin{itemize}
  \item In a commutative ring:

\[
(a)=\{r \cdot a \mid r \in R\}
\]

It is the set of all left (or right) multiples of \(a\) by elements from \(R\).
\item In a general (non-commutative) ring:
\begin{itemize}
\item Left principal ideal: \(\{r \cdot a \mid r \in R\}\)
\item Right principal ideal: \(\{a \cdot r \mid r \in R\}\)
\end{itemize}
\end{itemize}
\end{definition}

\begin{definition}[Principle Ideal Domain (PID)]

An integral domain in which all ideals are principal is called a Principal Ideal Domain (PID).
  
\end{definition}

\begin{theorem}
  The rings \(\mathbb{Z}\) and \(F[x], F\) a field, are PIDs. More generally, any Euclidean domain (ED) is a PID.
\end{theorem}

\begin{proofidea}
  
  
\end{proofidea}
  
\begin{formalproof}
  
\end{formalproof} 

\begin{theorem}
Any PID is a UFD
\end{theorem}

\section{Fields as Quotients}

\begin{theorem}[Quotients defining Fields]
Let \(F\) be a field and \(f(x) \in F[x]\) be irreducible. Then \(F[x] /(f(x))\) is a field, and is a vector space over \(F\) with basis \(B:=\left\{1, \bar{x}, \bar{x}^2, \ldots, \bar{x}^{n-1}\right\}\), where \(n=\operatorname{deg} f\). That is, every element in \(F[x] /(f(x))\) can be uniquely written as \(\overline{a_0+a_1 x+\cdots+a_{n-1} x^{n-1}}\), for \(a_i \in F\).
\end{theorem}

\begin{proofidea}
  
  
\end{proofidea}
  
\begin{formalproof}
  
\end{formalproof} 

\section{Finite Fields}

\begin{theorem}[Finite Fields]
For every prime \(p \in \mathbb{Z}\) and \(n \in \mathbb{N}\), there exists an irreducible polynomial \(f(x) \in\) \((\mathbb{Z} / p)[x]\) of degree \(n\). Thus \((\mathbb{Z} / p)[x] /(f(x))\) is a field with \(p^n\) elements. Any two such fields are isomorphic, and we denote the unique (up to isomorphism) field with \(p^n\) elements by \(\mathbb{F}_{p^n}\).
\end{theorem}


